% --- File: sections/03_methodology.tex ---

The core objective of our architecture is to dynamically forecast the 1-day ahead 99\% Value-at-Risk (VaR), denoted as $\text{VaR}_{t+1}^{0.01}$, utilizing historical high-frequency market data. To contextualize the hybrid approach, we first formalize the baseline models that represent the two ends of the modeling spectrum: the non-parametric historical approach and the unconstrained neural network.

\subsection{The Historical VaR Baseline}
Historical Simulation (HS) is the industry-standard non-parametric baseline, dominant in institutional risk systems due to its distributional agnosticism. Given a rolling window of $W$ past realized log-returns $\{r_{t-W+1}, \ldots, r_t\}$, the 1-day ahead VaR at confidence level $1 - q$ is simply the empirical quantile of that window:
\begin{equation}
    \text{VaR}_{t+1}^q = Q_q\!\left(r_{t-W+1}, \ldots, r_t\right)
\end{equation}
where $Q_q(\cdot)$ denotes the $q$-th order statistic of the sample. While parameter-free and straightforward to implement, Historical VaR suffers from the well-documented ``ghost effect'': a stress event influences the VaR estimate for exactly $W$ days, then drops out of the window abruptly, causing sudden artificial jumps. During volatility regime changes, this lag creates either persistent over-estimation (capital drag during calm regimes) or dangerous under-estimation (insufficient reserves immediately after a spike).

\subsection{The Naive LSTM Architecture}
To capture non-linear temporal dependencies and volatility clustering without strict distributional assumptions, we introduce a pure Deep Learning baseline using a Long Short-Term Memory (LSTM) network \cite{hochreiter1997long}. Let $\bm{x}_t$ be a feature vector containing historical returns, realized volatility, and trading volumes. The LSTM maintains a hidden state $\bm{h}_t$ and a cell state $\bm{c}_t$, updated as follows:
\begin{equation}
    \bm{h}_t, \bm{c}_t = \text{LSTM}(\bm{x}_t, \bm{h}_{t-1}, \bm{c}_{t-1}; \bm{\Theta})
\end{equation}
where $\bm{\Theta}$ represents the learned network parameters. The output layer projects the hidden state to a VaR forecast: $\widehat{\text{VaR}}_{t+1}^q = \bm{W}_{out} \bm{h}_t + b_{out}$.

However, training a ``Naive'' LSTM using standard symmetric loss functions like Mean Squared Error (MSE) is fundamentally flawed for risk applications. MSE optimizes for the conditional mean, whereas VaR is strictly a quantile property. Furthermore, an unconstrained neural network operating in a noisy, non-stationary environment---such as our extreme stress-testing proxy---is highly susceptible to overfitting, learning market noise rather than true structural risk dynamics.

\subsection{Anchored Quantile Regression (Physics-Informed)}
To combine the interpretability of the Historical VaR baseline with the adaptive capacity of the LSTM, we introduce an ``Anchored'' architecture. We achieve this by redefining the network's objective function using an Asymmetric Quantile Loss (also known as the Pinball Loss) \cite{koenker1978regression} and injecting a rolling Parametric VaR estimate as a structural prior---akin to Physics-Informed Neural Networks (PINNs) where physical laws constrain the solution space. This separation is deliberate: the Historical VaR serves as the empirical comparison benchmark, while the parametric prior acts as the safety regularizer inside the loss function.

The network is trained to directly minimize the asymmetric loss $L_q$ for a given target quantile $q$:
\begin{equation}
    L_q(r_{t+1}, \widehat{\text{VaR}}_{t+1}^q) = 
    \begin{cases} 
      q (r_{t+1} - \widehat{\text{VaR}}_{t+1}^q) & \text{if } r_{t+1} \geq \widehat{\text{VaR}}_{t+1}^q \\
      (1 - q) (\widehat{\text{VaR}}_{t+1}^q - r_{t+1}) & \text{if } r_{t+1} < \widehat{\text{VaR}}_{t+1}^q
    \end{cases}
\end{equation}

In our framework, the total loss function $\mathcal{L}_{total}$ is augmented with a regularization penalty $\lambda$ that ``anchors'' the neural network's output to the rolling Parametric VaR estimate ($\text{VaR}_{param}$):
\begin{equation}
    \mathcal{L}_{total} = \frac{1}{N} \sum_{t=1}^{N} L_q(r_{t+1}, \widehat{\text{VaR}}_{t+1}^q) + \lambda \left\| \widehat{\text{VaR}}_{t+1}^q - \text{VaR}_{param, t+1}^q \right\|^2
\end{equation}

This architectural design ensures that the deep learning model focuses on optimizing the 1st percentile of the return distribution (via the Pinball Loss) while the anchoring penalty $\lambda$ acts as a regulatory safety rail. It prevents the LSTM from generating nonsensical predictions during unprecedented Black Swan events, effectively resolving the "black box" vulnerability and ensuring compliance with Basel IV backtesting standards.